\section{The exact phase-graph map to the original averaged quotient and its completion}
\label{sec:pgb-averaged-bridge}

Fix the original full quartet and all the data of PGS.1--PGS.41.
Thus $m,k\geq1$, $d=4m$, $\rho=1/2+\delta+i\gamma$,
$0<\delta<1/2$, $\gamma>2$,
\[
 h(s)=\prod_{r\in\{\rho,\bar\rho,1-\rho,1-\bar\rho\}}(s-r)^m,
 \qquad v_h(s)=\frac{2\xi(s)}{h(s)},\qquad
 \Phi'=h,\quad\Phi(0)=0.
\]
Retain $s_i=1/2+it_i$, $\Sigma=\sum_i s_i=k/2+iu$,
$u=\sum_i t_i$, $\Phi(1/2+it)=A+iq_\Phi(t)$, $A>0$,
and $\tau_k=\sum_i\Phi(s_i)=kA+i\sum_iq_\Phi(t_i)$.
The polynomial $q_\Phi$ has degree $4m+1$ and leading
coefficient $1/(4m+1)$. These are the original objects, and the
full-divisor hypothesis of the original counterfactual packet is
unchanged. Let
\[
 \begin{gathered}
 \mathscr H_k=L^2\left(\mathbb{R}^k,\prod_i\frac{dt_i}{2\pi}\right),
 \qquad \mathsf T_k f=\tau_k f,
 \quad\operatorname{dom}\mathsf T_k
 =\{f\in\mathscr H_k:\tau_kf\in\mathscr H_k\},\\
 \chi=\chi_{h,k},\quad q=[1+k(m-1)](k+1)^2,
 \quad E=\mathbb{C}[S]/(\chi),\\
 \mathcal P_N=\mathbb{C}[S]_{\leq N},\quad N\geq q-1,
 \qquad B_N:\mathcal P_{N-q}\longrightarrow\mathcal P_N,
 \quad Q\longmapsto\chi Q,\\
 J_N:\mathcal P_N\longrightarrow E,\quad P\longmapsto[P],
 \qquad s:E\longrightarrow\mathcal P_{q-1},
 \quad [P]\longmapsto\operatorname{rem}_{\chi}P.
 \end{gathered}
 \tag{PGB.1}\label{eq:pgb-data}
\]
Here $\mathcal P_{-1}=0$. Every finite matrix below uses the
original increasing monomial frame, including the increasing
remainder frame of $E$. All Hilbert inner products are conjugate
linear in their first argument. Monic division proves
$\ker J_N=\operatorname{im}B_N$, injectivity of $B_N$, and
$J_Ns=I_E$.

Write the actual polynomial observation and defect as
\[
 \begin{gathered}
 (\mathcal V_NP)(t)=\Bigl(\prod_i v_h(s_i)\Bigr)P(\Sigma),
 \qquad Lx=\Bigl(\prod_i v_h(s_i)\Bigr)F_x(s_1,\ldots,s_k),\\
 F_x=\sum_{b=0}^{q-1}x_b\,
 \operatorname{rem}_{h(s_1),\ldots,h(s_k)}(\tau_k\Sigma^b),
 \qquad D_N=\mathsf T_k\mathcal V_N-LJ_N,\\
 \mathcal A_N=(\mathcal V_N,D_N):
 \mathcal P_N\longrightarrow\mathscr H_k\oplus\mathscr H_k,
 \qquad M_N=\mathcal V_N^*\mathcal V_N,
 \quad\widehat M_N=\mathcal A_N^*\mathcal A_N=M_N+D_N^*D_N.
 \end{gathered}
 \tag{PGB.2}\label{eq:pgb-original-maps}
\]
The definition of $F_x$ retains every original ordered tensor
remainder and every factor $v_h$. Its equality with
$\mathsf S_kC_k\eta x$ and the injectivity of $L$ are the
fully proved PGS.4 construction. In particular $F_x=0$ implies
$x=0$. The source decay of PGS and PGD places all the displayed
columns in the multiplication domain. The injectivity of
$\mathcal V_N$ makes both finite Grams positive definite.

\subsection{The positive cost of adding the literal defect}

Suppress $N$ only inside the next calculation. Define
\[
 \begin{gathered}
 H=B^*MB,\qquad
 C=s-BH^{-1}B^*Ms,\qquad G=C^*MC=(JM^{-1}J^*)^{-1},\\
 T=DB:\mathcal P_{N-q}\longrightarrow\mathscr H_k,
 \qquad H_+=H+T^*T=B^*\widehat M B,\\
 Q_D=I_{\mathscr H_k}-TH_+^{-1}T^*.
 \end{gathered}
 \tag{PGB.3}\label{eq:pgb-augmentation-data}
\]
If $N=q-1$, every map through the zero relation space is zero,
both relation determinants are one, and $Q_D=I$.
For positive relation dimension, $H,H_+$ are positive definite.
The formula for $C$ is proved by writing a lift as $sx+By$ and
completing its $M$ square. It gives $JC=I$ and $B^*MC=0$.
To verify the inverse formula for $G$, decompose every polynomial
as $CJp+By$; pairing against $Cx$ proves $MC=J^*G$.
Multiplication first by $M^{-1}$ and then by $J$ gives the formula.

The bounded operator $Q_D$ is positive and has the exact inverse
\[
 Q_D=(I_{\mathscr H_k}+TH^{-1}T^*)^{-1},\qquad
 0<Q_D\preceq I_{\mathscr H_k}.
 \tag{PGB.4}\label{eq:pgb-positive-inverse}
\]
Indeed $I+TH^{-1}T^*$ is self-adjoint and at least $I$.
The product of this operator with $I-TH_+^{-1}T^*$ is $I$:
the coefficient between $T$ and $T^*$ in the remaining term is
\[
 H^{-1}-H_+^{-1}-H^{-1}T^*TH_+^{-1}
 =H^{-1}-(I+H^{-1}T^*T)H_+^{-1}=0.
\]
Taking adjoints proves the product in the opposite order.
On the finite-dimensional range of $T$ positivity follows by
diagonalizing the positive operator $TH^{-1}T^*$; on its
orthogonal complement the inverse is $I$. This also proves
$Q_D\succeq(1+\|TH^{-1}T^*\|)^{-1}I$ and its bounded
positive square root without requiring an infinite-dimensional
spectral assertion.

Every lift is $Cx+By$. Its unchanged graph norm is
\[
 \begin{aligned}
 \|\mathcal A(Cx+By)\|^2
 &=x^*Gx+y^*Hy+\|DCx+Ty\|^2\\
 &=x^*\bigl(G+C^*D^*Q_DDC\bigr)x
 +(y+H_+^{-1}T^*DCx)^*H_+
                  (y+H_+^{-1}T^*DCx).
 \end{aligned}
 \tag{PGB.5}\label{eq:pgb-augmentation-square}
\]
Expansion proves this identity with its two positive cross terms
in the first square and the negative minimized contribution
$-C^*D^*TH_+^{-1}T^*DC$ retained. Therefore
\[
 \begin{gathered}
 \widehat R_N=C_N-B_N(H_{+,N})^{-1}T_N^*D_NC_N,\\
 \widehat G_N=G_N+\mathcal E_N^{\rm aug},\qquad
 \mathcal E_N^{\rm aug}=C_N^*D_N^*Q_{D,N}D_NC_N>0,\\
 \ker\mathcal E_N^{\rm aug}=\ker(D_NC_N)=\{0\}.
 \end{gathered}
 \tag{PGB.6}\label{eq:pgb-augmentation-quotient}
\]
The section has $J_N\widehat R_N=I$ and is orthogonal to the
relation columns in $\widehat M_N$. Thus uniqueness of a minimum
identifies it with the original section in PGS.13 and PGS.20;
the displayed $\widehat G_N$ is that same graph quotient.
The kernel assertion follows from the strict bounded positivity
of $Q_{D,N}$ proved above. The actual map whose Gram is the
new positive term is
$x\mapsto Q_{D,N}^{1/2}D_NC_Nx$ into $\mathscr H_k$.

Here the kernel can be completely evaluated on the actual source.
If $P\ne0$ has degree $n$ and leading coefficient $p_n$, the
polynomial $\tau_kP(\Sigma)$, as a polynomial in $s_1$, has
degree $n+d+1$ and leading coefficient $p_n/(d+1)$.
Indeed the original monic $h$ gives $\Phi$ leading coefficient
$1/(d+1)$, while the terms $\Phi(s_j)P(\Sigma)$ with $j>1$
have $s_1$ degree only $n$. On the other hand every variable
degree of $F_{JP}$ is at most $d-1$ by its original successive
remainders. Consequently $\tau_kP(\Sigma)-F_{JP}$ is a nonzero
polynomial. Its restriction under $s_i=1/2+it_i$ is nonzero
almost everywhere outside its polynomial zero set; the product
$\prod_i v_h(s_i)$ is nonzero almost everywhere because its
one-variable factors are nonzero entire functions. Thus
$D_NP\ne0$ in $\mathscr H_k$. This proves injectivity of
$D_N$. Since $J_NC_N=I$, it also proves injectivity of
$D_NC_N$, all asserted zero kernels, and the strict sign in
PGB.6. In particular $G_N<\widehat G_N$ even at the empty
relation window.

There is a second realization retaining both original graph
components. Put $\mathfrak a_N=\widehat R_N-C_N:E\to\operatorname{im}B_N$.
Orthogonality of $\mathcal V_NC_N$ and $\mathcal V_N\mathfrak a_N$ gives
\[
 \begin{gathered}
 \mathcal A_N\widehat R_N
 =(\mathcal V_NC_N,0)
       +(\mathcal V_N\mathfrak a_N,D_N\widehat R_N),\\
 \mathcal E_N^{\rm aug}=\mathfrak a_N^*M_N\mathfrak a_N
        +(D_N\widehat R_N)^*(D_N\widehat R_N).
 \end{gathered}
 \tag{PGB.7}\label{eq:pgb-augmentation-map}
\]
The two displayed summands in the first line are orthogonal.
Taking their Grams proves the second line and also proves its
equality with PGB.6, rather than merely comparing two presentations.

\subsection{Composition with the accepted all-holonomy quotient}

At each fixed $N$, retain the original HCD circumference, denoted
here by $L_{\rm circ}$ to distinguish that scalar from the
observation map $L:E\to\mathscr H_k$. Thus $L_{\rm circ}$
is exactly the quantity called $L$ in HCD, with the same value.
Retain the original holonomy matrices of HCD.2--HCD.2c.
Concretely choose the already
computed original $\beta>0$, weighted source constant
$\kappa_{\beta,N}$ and any $0<\eta<1$, and take
$L_{\rm circ}\geq\beta^{-1}\log(1+\kappa_{\beta,N}/\eta)$.
With $d\mu=d\theta/(2\pi)$ the HCD proof then gives
\[
 \begin{gathered}
 \int_0^{2\pi}M_{N,\theta}\,d\mu=M_N,
 \qquad(1-\eta)M_N\preceq M_{N,\theta}
                              \preceq(1+\eta)M_N,\\
 C_{N,\theta}=M_{N,\theta}^{-1}J_N^*G_{N,\theta},\qquad
 G_{N,\theta}=(J_NM_{N,\theta}^{-1}J_N^*)^{-1},\qquad
 \overline G_N=\int_0^{2\pi}G_{N,\theta}\,d\mu.
 \end{gathered}
 \tag{PGB.8}\label{eq:pgb-accepted-holonomy}
\]
No phase graph matrix is substituted into this accepted definition.
The full relative Hilbert fibre and the original half-density
source used to construct $M_{N,\theta}$ remain those of HCD.
Uniform finite-dimensional equivalence makes every following
coefficient field measurable and square integrable.

For each phase set
\[
 X_{N,\theta}=(B_N^*M_{N,\theta}B_N)^{-1}
                          B_N^*M_{N,\theta}C_N,
 \qquad C_N-C_{N,\theta}=B_NX_{N,\theta}.
 \tag{PGB.9}\label{eq:pgb-original-relation-field}
\]
At zero relation dimension $X_{N,\theta}$ is the zero map.
Both sections lift the same vector, so their difference lies in
$\operatorname{im}B_N$. Pairing the difference with $B_N$ in
the phase metric, and using $B_N^*M_{N,\theta}C_{N,\theta}=0$,
proves exactly the displayed formula. Squaring this orthogonal
decomposition of $C_N$ and integrating proves
\[
 \mathscr D_N^{\rm hol}:=G_N-\overline G_N
 =\int_0^{2\pi}(C_N-C_{N,\theta})^*M_{N,\theta}
                      (C_N-C_{N,\theta})\,d\mu\succeq0.
 \tag{PGB.10}\label{eq:pgb-original-gap}
\]
This is the exact accepted HCD.13 relation cost. In particular
the graph augmentation does not remove it.

Define
\[
 \mathscr B_N=L^2(\Theta;\operatorname{im}B_N,
                                  M_{N,\theta}\,d\mu),
 \quad
 \mathfrak F_Nx=
 \left(Q_{D,N}^{1/2}D_NC_Nx,
       [\theta\mapsto(C_N-C_{N,\theta})x]\right)
 \in\mathscr H_k\oplus\mathscr B_N.
 \tag{PGB.11}\label{eq:pgb-full-defect-map}
\]
Its complete matrix identity and kernel are
\[
 \begin{gathered}
 \boxed{\quad
 \widehat G_N-\overline G_N
 =\mathcal E_N^{\rm aug}+\mathscr D_N^{\rm hol}
 =\mathfrak F_N^*\mathfrak F_N>0,\quad}\\
 \ker\mathfrak F_N
 =\{x:D_NC_Nx=0,\ C_Nx=C_{N,\theta}x\text{ a.e.}\}=\{0\}.
 \end{gathered}
 \tag{PGB.12}\label{eq:pgb-main-bridge}
\]
Indeed the two orthogonal direct-sum target norms are precisely
PGB.6 and PGB.10. A sum of two nonnegative real norms vanishes
exactly when both vanish; the positivity in PGB.4 and that of
the phase source metric give the stated full kernel.
The actual injectivity proved after PGB.6 evaluates that kernel
as zero and makes the displayed matrix difference strictly positive.

The identity on coefficients consequently defines the actual
contraction $(E,\widehat G_N)\to(E,\overline G_N)$. It is
induced by an explicit map of the original two-term complexes:
send $p$ in $(\mathcal P_N,\widehat M_N)$ to the constant
field $\theta\mapsto p$ in the HCD common-image space, and
send a relation $B_Ny$ to that same constant relation field.
The inclusion differentials commute, $J_Np$ is unchanged, and
\[
 \int p^*M_{N,\theta}p\,d\mu=p^*M_Np
                  \leq p^*\widehat M_Np.
\]
Thus the source map is a contraction and its induced quotient
map is the displayed identity. More precisely, the complete
isometric dilation of this quotient contraction is
\[
 x\longmapsto
 \bigl(x,Q_{D,N}^{1/2}D_NC_Nx,
             [\theta\mapsto(C_N-C_{N,\theta})x]\bigr)
 \quad\hbox{from }(E,\widehat G_N)
 \hbox{ to }(E,\overline G_N)\oplus\mathscr H_k\oplus\mathscr B_N.
 \tag{PGB.13}\label{eq:pgb-isometric-dilation}
\]
Its squared norm identity is PGB.12; polarization proves every
cross pairing. This gives the precise relation between the graph
quotient, the accepted average, and the common finite-source
relation obstruction on the unchanged coefficient object $E$.

\subsection{Exponential moments and the complete graph quotient}

For clarity the density argument needed for the degree limit is
proved here. If a finite positive Borel measure $\nu$ on
$\mathbb{R}$ has $\int e^{\epsilon|u|}\,d\nu<\infty$ for some
$\epsilon>0$, then polynomials are dense in $L^2(\nu)$.
To prove this, let $f$ be orthogonal to them and put
$d\lambda=\bar f\,d\nu$. Cauchy--Schwarz gives
\[
 \int e^{\epsilon|u|/2}\,|d\lambda|
 \leq\|f\|_{L^2(\nu)}
       \left(\int e^{\epsilon|u|}\,d\nu\right)^{1/2}<\infty.
\]
The Fourier integral $F(z)=\int e^{izu}\,d\lambda(u)$ is
holomorphic on $|\operatorname{Im}z|<\epsilon/2$:
on a smaller closed strip every differentiated integrand is
dominated by a constant times $e^{\epsilon|u|/2}|d\lambda|$.
Every derivative at zero vanishes by polynomial orthogonality.
The Taylor expansion and then overlapping analytic discs give
$F=0$ throughout this connected strip, in particular on its
real axis.

Here is the needed uniqueness on that axis. Let
$g_t(x)=(4\pi t)^{-1/2}\exp(-x^2/(4t))$, $t>0$.
The elementary Gaussian integral gives
\[
 g_t(v)=\frac1{2\pi}\int_{\mathbb{R}}e^{-t\zeta^2}e^{i\zeta v}
                                                  \,d\zeta.
\]
For completeness, differentiating the latter integral in $v$
and integrating by parts in $\zeta$ gives the differential
equation $I'(v)=-vI(v)/(2t)$; its value at zero is
$I(0)=\sqrt{\pi/t}$, obtained by squaring the Gaussian integral
and using polar coordinates. These statements prove the formula
with its constants. Absolute integrability and Fubini now give
$\lambda*g_t=0$ because $F(-\zeta)=0$.
For $\varphi\in C_c(\mathbb{R})$ it follows that
$\int(\varphi*g_t)\,d\lambda=0$. The functions
$\varphi*g_t$ converge uniformly to $\varphi$: split the
unit-mass Gaussian integral into $|y|<a$, where uniform
continuity controls the difference, and $|y|\geq a$, where
its mass tends to zero. Thus $\int\varphi\,d\lambda=0$ for
every such $\varphi$. Regularity of a finite Borel measure on
the real line, followed by continuous approximations to indicators
of compact sets and open sets, gives $\lambda=0$.
It follows that $f=0$ $\nu$-almost everywhere. The orthogonal
complement criterion for closed subspaces proves the density claim.

Apply this to the original two densities
\[
 \begin{gathered}
 w(t)=|v_h(1/2+it)|^2/(2\pi),\quad m_k=w^{*k},\\
 r_k(u)=\int_{\sum t_i=u}(1+|\tau_k|^2)\prod_iw(t_i),
 \qquad\mathscr L_k=L^2(\mathbb{R},r_k(u)\,du).
 \end{gathered}
 \tag{PGB.14}\label{eq:pgb-original-density}
\]
All fibre integrals use elimination of $t_k$ and measure
$dt_1\cdots dt_{k-1}$, with determinant-one change of variables.
For $k=1$ the fibre integral is evaluation.
The proved PGD.11 bound is
$w(t)\leq C_he^{-\alpha|t|}(1+|t|)^{-p_0}$,
$\alpha=\pi/2$, $p_0=8m-9/2>1$.
For fixed $k,h$ and $0<\epsilon<\alpha$, both
\[
 \int e^{\epsilon|u|}r_k(u)\,du<\infty,
 \qquad
 \int e^{\epsilon|u|}|\chi(k/2+iu)|^2r_k(u)\,du<\infty
 \tag{PGB.15}\label{eq:pgb-actual-exponential-moments}
\]
follow directly on the product source. Indeed
$|u|\leq\sum_i|t_i|$, while $1+|\tau_k|^2$ and
$|\chi(k/2+iu)|^2$ are bounded by fixed products of
powers of $(1+|t_i|)$. Their product with
$\prod_i e^{-(\alpha-\epsilon)|t_i|}$ is integrable.
The same proof holds with $r_k$ replaced by $m_k$.
Every constant here depends on the fixed original packet and
$k$; no uniform statement in $k$ has been inserted.

The density is positive almost everywhere by PGS.35, and the
polynomial $\chi(k/2+iu)$ has only finitely many real zeros.
Consequently multiplication by that unchanged polynomial is
the onto isometry
\[
 L^2(\mathbb{R},|\chi(k/2+iu)|^2r_k(u)\,du)
 \longrightarrow\mathscr L_k,
 \qquad f(u)\longmapsto\chi(k/2+iu)f(u).
 \tag{PGB.16}\label{eq:pgb-relation-density-map}
\]
Its inverse is division by $\chi(k/2+iu)$ away from its finite
zero set, with arbitrary values on that null set. Equality of
the two squared norms proves both the isometry and the inverse
domain assertion. PGB.15 and the density proof therefore imply
\[
 \overline{\{\chi(k/2+iu)Q(k/2+iu):Q\in\mathbb{C}[S]\}}
              ^{\,\mathscr L_k}=\mathscr L_k.
 \tag{PGB.17}\label{eq:pgb-dense-actual-relations}
\]
Polynomials in $u$ and in $S=k/2+iu$ have exactly the same span:
the forward substitution has inverse $u=(S-k/2)/i$.
This proves the density assertion on the original coordinate,
not on an independently substituted source.

Recall the actual isometry and observation from PGS.36--PGS.40:
\[
 \begin{gathered}
 \mathscr J_k:\mathscr L_k\longrightarrow
                      \mathscr H_k\oplus\mathscr H_k,
 \qquad
 \mathscr J_k f=
 \left(\prod_iv_h(s_i)f(u),\tau_k\prod_iv_h(s_i)f(u)\right),\\
 Y=(0,L),\qquad g=\mathscr J_k^*Y,
 \quad g_x(u)=z(u)x/r_k(u),\quad
 z_b(u)=\int_{\sum t_i=u}\overline{\tau_k}F_b(s)\prod_iw(t_i),\\
 Y_\perp=(I-\mathscr J_k\mathscr J_k^*)Y,
 \qquad\Omega_k=Y_\perp^*Y_\perp>0.
 \end{gathered}
 \tag{PGB.18}\label{eq:pgb-full-observation}
\]
These maps are the original isometry and mixed observation;
the adjoint is exactly $z/r_k$ with its conjugate phase.
To recall why the strict sign holds, fibre Cauchy--Schwarz gives
$|z(u)x|^2\leq(r_k(u)-m_k(u))x^*\Lambda(u)x$, where
$\Lambda_{bc}(u)=\int_{\sum t_i=u}\overline{F_b}F_c\prod_iw$.
It follows that
$\Omega_k\succeq\int(m_k/r_k)\Lambda\,du>0$;
for nonzero $x$, the integral of $x^*\Lambda x$ is
$\|Lx\|^2>0$, and $m_k/r_k>0$ almost everywhere.
This argument also proves $g_x\in\mathscr L_k$ and the stated
adjoint by direct pairing. Thus the preceding uses of these
objects involve actual bounded maps on the specified domains.

Put $\iota P(u)=P(k/2+iu)$. For every original polynomial $P$
write $\mathcal A P$ and $JP$ for the compatible union of the
finite-window maps, whose definitions agree on every inclusion.
Then
the exact source identity is
\[
 \mathcal A P=
 \mathscr J_k(\iota P-gJ P)-Y_\perp J P,
 \qquad
 \|\mathcal A P\|^2=
 \|\iota P-gJ P\|_{\mathscr L_k}^2+(JP)^*\Omega_k(JP).
 \tag{PGB.19}\label{eq:pgb-completion-source-identity}
\]
It follows by substituting
$Y=\mathscr J_kg+Y_\perp$ into $\mathcal A=U-YJ$.
The two summands lie in orthogonal closed subspaces. Define
\[
 \begin{gathered}
 \mathscr K_\infty=
 \mathscr J_k\mathscr L_k\ \mathbin{\oplus^\perp}\ Y_\perp E,
 \qquad
 \Xi:\mathscr L_k\oplus(E,\Omega_k)\longrightarrow\mathscr K_\infty,
 \quad (f,x)\longmapsto\mathscr J_k f-Y_\perp x,\\
 \Xi^{-1}(v)=
 \left(\mathscr J_k^*v,-\Omega_k^{-1}Y_\perp^*v\right).
 \end{gathered}
 \tag{PGB.20}\label{eq:pgb-completion-isometry}
\]
The space is closed because the first summand is the range of
an isometry and the second is finite dimensional and orthogonal.
The formulas $\mathscr J_k^*Y_\perp=0$ and
$Y_\perp^*Y_\perp=\Omega_k$ prove both inverse identities,
and prove the exact squared norm
$\|\Xi(f,x)\|^2=\|f\|^2+x^*\Omega_kx$.
This is a complete isometric parametrization with its minus
sign determined by the original subtraction in $D$.

In fact $\mathscr K_\infty$ is exactly the closure of the
original polynomial graph image. Inclusion follows from PGB.19.
Conversely take $(f,x)$ in the domain of $\Xi$.
The affine set
\[
 \{\iota(sx+\chi Q)-g_x:Q\in\mathbb{C}[S]\}
\]
is dense in $\mathscr L_k$ by PGB.17. Choose $Q_j$ for which
these functions converge to $f$ and put $P_j=sx+\chi Q_j$.
Then $JP_j=x$ for every $j$, and PGB.19 proves
$\mathcal AP_j\to\Xi(f,x)$. This proves surjectivity onto
the completed graph from limits of original polynomials,
retaining the prescribed quotient coordinate throughout.
In the same manner
\[
 \overline{\mathcal A(\chi\mathbb{C}[S])}
      =\mathscr J_k\mathscr L_k,
 \qquad
 \mathscr K_\infty/\overline{\mathcal A(\chi\mathbb{C}[S])}
       \simeq(E,\Omega_k).
 \tag{PGB.21}\label{eq:pgb-completed-quotient}
\]
The quotient map on the whole completed source is the bounded
map $v\mapsto-\Omega_k^{-1}Y_\perp^*v$; it sends
$\mathcal AP$ to $JP$ and has precisely the displayed closed
relation space as kernel. Its section is $x\mapsto-Y_\perp x$.
This explicitly computes the cochain completion and every map.

This completed source is a closed linear relation in
$\mathscr H_k\oplus\mathscr H_k$, with its vertical part fully
determined. Indeed $Y=\mathscr J_kg+Y_\perp$ gives
\[
 \Xi(f,x)=\mathscr J_k(f+g_x)-Yx,
 \qquad
 \mathscr K_\infty\cap(\{0\}\oplus\mathscr H_k)
          =\{(0,-Lx):x\in E\}.
 \tag{PGB.21a}\label{eq:pgb-vertical-kernel}
\]
If the first component vanishes, its squared norm is
$\int|f+g_x|^2m_k\,du=0$. Since $m_k>0$ almost everywhere,
$f=-g_x$ almost everywhere, also as elements of $\mathscr L_k$;
the source vector is therefore $-Yx$. Conversely
$\Xi(-g_x,x)=-Yx$ proves the reverse inclusion.
In particular the original densely defined defect operator
$\mathcal VP\mapsto DP$ on the closure of the polynomial
observation source is not closable. For any fixed $x\ne0$,
PGB.17 supplies $P_j=sx+\chi Q_j$ with
$\iota P_j\to0$ in $\mathscr L_k$. The original first
components tend to zero, while
$DP_j=\mathsf T_k\mathcal VP_j-Lx\to-Lx\ne0$.
The two convergence assertions follow from the two component
norms of $\mathscr J_k\iota P_j$. This witnesses the entire
vertical subspace without asserting that the completed relation
is the graph of a single-valued operator.

There is also an explicit map to the accepted unweighted
completion. Put $\mathscr L_k^0=L^2(m_kdu)$ and
\[
 \mathscr J_k^0:\mathscr L_k^0\to\mathscr H_k,
 \quad f\mapsto\prod_i v_h(s_i)f(u),\qquad
 I_{r,m}:\mathscr L_k\to\mathscr L_k^0,\quad f\mapsto f.
\]
The first is an isometry by the original convolution integral;
the second is a contraction because $r_k\geq m_k$.
The first-coordinate projection $\pi_1$ obeys
$\pi_1\Xi(f,x)=\mathscr J_k^0I_{r,m}(f+g_x)$ and
$\pi_1\mathscr J_k f=\mathscr J_k^0I_{r,m}f$.
Since the same exponential-moment argument proves that the
unweighted relation closure is all of $\mathscr J_k^0\mathscr L_k^0$,
these formulas define a bounded cochain map
\[
 [\mathscr J_k\mathscr L_k\hookrightarrow\mathscr K_\infty]
 \xrightarrow{\ \pi_1\ }
 [\mathscr J_k^0\mathscr L_k^0\xrightarrow{I}
                         \mathscr J_k^0\mathscr L_k^0].
 \tag{PGB.21b}\label{eq:pgb-to-unweighted-completion}
\]
It agrees on the original polynomial complex with
$\mathcal AP\mapsto\mathcal VP$ and
$\mathcal A(\chi Q)\mapsto\mathcal V(\chi Q)$.
Its degree-one map is $(E,\Omega_k)\to0$ with kernel all
of $E$. Thus the positive completed graph quotient maps exactly
to the zero unweighted completion in the accepted R66 source;
its nonzero metric and the receiving zero object are connected
by this specific cochain map and its full kernel.

\subsection{The degree limit and the exact obstruction to a uniform comparison}

For any $x\in E$, minimizing PGB.19 over $P\in\mathcal P_N$
with $JP=x$ gives
\[
 x^*\widehat G_Nx
 =x^*\Omega_kx+
 \inf_{Q\in\mathcal P_{N-q}}
       \|\iota(sx+\chi Q)-g_x\|_{\mathscr L_k}^2.
 \tag{PGB.22}\label{eq:pgb-quotient-distance}
\]
At $N=q-1$ the infimum contains exactly $Q=0$.
The nested relation spaces and their proved density show that
the nonnegative second term decreases to zero for every $x$.
Using polarization on the fixed finite remainder frame gives
entrywise convergence, hence operator-norm convergence of the
finite matrices. Thus, at fixed $h,k$,
\[
 \widehat G_N\downarrow\Omega_k>0,
 \qquad G_{W,N}\downarrow0,
 \qquad\Delta_N\downarrow\Omega_k,
 \qquad G_N\downarrow0.
 \tag{PGB.23}\label{eq:pgb-three-limits}
\]
For $G_{W,N}$ repeat the distance calculation with $g=0$ and
without the $\Omega_k$ term, using the same $r_k$ and PGB.17.
For $\Delta_N$, its PGS.40 expression is
$\Omega_k+((I-\mathsf P_N)g)^*((I-\mathsf P_N)g)$;
polynomial density in $\mathscr L_k$ makes the latter term
decrease to zero. For $G_N$ repeat the affine approximation
argument in $L^2(m_kdu)$, whose exponential moments and dense
relation image were proved by the same PGB.15--PGB.17 argument.
These are matrix limits in the unchanged frame, not estimates
in the separate tensor variable $k$.

The original averaged minimum form is defined at every
$L_{\rm circ}>0$, including a phase with a singular finite
matrix, by its actual source map $T_\theta$ into the sampled
Hilbert space of HCD.18. Let $P_\theta^{\rm rel}$ be orthogonal
projection onto the finite-dimensional closed subspace
$T_\theta(\operatorname{im}B_N)$. Then
\[
 \begin{aligned}
 x^*G_{N,\theta}x
 &=\|(I-P_\theta^{\rm rel})T_\theta C_Nx\|^2,\\
 x^*(G_N-\overline G_N)x
 &=\int\|P_\theta^{\rm rel}T_\theta C_Nx\|^2\,d\mu.
 \end{aligned}
\]
The first equality follows by minimizing over the actual affine
source image; the second follows from its orthogonal
decomposition and $\int\|T_\theta C_Nx\|^2d\mu=x^*G_Nx$.
The projections are measurable: for the map $T_\theta B_N$,
their finite-rank formula is the pointwise limit as $t\downarrow0$
of $(T_\theta B_N)(B_N^*M_{N,\theta}B_N+tI)^{-1}
(T_\theta B_N)^*$. Diagonalization of the finite Gram proves
this limit, including its zero eigenvalues. The displayed
integrands are bounded above by the integrable full source
norm, so both minimum forms and integrals are defined.
Consequently $0\preceq\overline G_N\preceq G_N$ at every
original circumference, with no inverse imposed on a singular
phase. Positivity of the finite $\overline G_N$ follows directly:
for $k\geq2$ every original sampled convolution weight is
positive; for $k=1$ all weights are positive outside a countable
set of phases determined by the isolated real zeros of $v_h$.
At such a phase a nonzero polynomial cannot vanish on infinitely
many distinct sample coordinates, so its finite source Gram
and its quotient Gram are positive definite. Integrating any
nonzero quotient vector gives a strictly positive norm.

The identity $G_N\to0$ is therefore precisely consistent with
the accepted R66/CR conclusion. For these actual averaged
minimum forms the augmentation identity gives
\[
 \mathcal E_N^{\rm aug}=\widehat G_N-G_N\longrightarrow\Omega_k,
 \qquad
 \widehat G_N-\overline G_N\longrightarrow\Omega_k
 \tag{PGB.24}\label{eq:pgb-comparison-limit}
\]
at every fixed $L_{\rm circ}>0$.
This conclusion also holds for any choices $L_{\rm circ}=L_N$ satisfying
PGB.8 at each finite cutoff; no interchange of the limit with
a changing phase source was used, since the order bound alone
gives $\|\overline G_N\|\leq\|G_N\|\to0$.

Where $\overline G_N>0$, which the accepted source proves at
every finite $N$, the comparison has the explicit divergence
\[
 \inf_{x\ne0}
   \frac{x^*\widehat G_Nx}{x^*\overline G_Nx}
 \ \geq\ \frac{\lambda_{\min}(\Omega_k)}{\|G_N\|}
 \ \longrightarrow\ +\infty.
 \tag{PGB.25}\label{eq:pgb-uniform-obstruction}
\]
Indeed the numerator is at least
$\lambda_{\min}(\Omega_k)\|x\|^2$, while the denominator
is at most $\|G_N\|\|x\|^2$. Both constants refer to the same
original remainder frame. Hence no finite constant independent
of $N$ can bound $\widehat G_N$ above by that constant times
the accepted $\overline G_N$. The direction
$\overline G_N\preceq\widehat G_N$ remains the exact
cochain-induced contraction PGB.12--PGB.13. This proves an
actual relation and its precise obstruction, rather than
inferring lack of relation from the obstruction.

Finally fix the same original circumference at all four endpoints.
For the inverse-form realization PGB.8, choose it using the
original common maximal source $N_*=2q+2$ as in R63 and restrict
its bound to every endpoint source. Alternatively use the
proved minimum-form realization at any fixed circumference.
The literal endpoint convention remains
$(q-1,q,2q-1,2q)$ with signs $(+,+,-,-)$. Set
$K_N^{\rm av}=\overline G_N^{-1/2}\widehat G_N
\overline G_N^{-1/2}>I$ at these finite endpoints.
Direct determinant multiplication gives the exact signed bridge
\[
 \begin{aligned}
 &\bigl[\log\det\widehat G_{q-1}+\log\det\widehat G_q
        -\log\det\widehat G_{2q-1}-\log\det\widehat G_{2q}\bigr]\\
 &\quad-
 \bigl[\log\det\overline G_{q-1}+\log\det\overline G_q
        -\log\det\overline G_{2q-1}-\log\det\overline G_{2q}\bigr]\\
 &=\log\det K_{q-1}^{\rm av}+\log\det K_q^{\rm av}
       -\log\det K_{2q-1}^{\rm av}-\log\det K_{2q}^{\rm av}.
 \end{aligned}
 \tag{PGB.26}\label{eq:pgb-four-endpoint-bridge}
\]
Each separate logarithm on the right is nonnegative; their
signed sum has no sign consequence from that fact alone.
The identity retains all four original endpoints, including
their possible repeated appearance when $q=1$.

All the maps on $E$ and polynomial coefficient spaces above
extend to a fixed original active mask by applying the same
linear map in each coordinate. For PGB.12 the represented-zero
kernel is exactly the stated kernel in each coordinate; for
the completed quotient PGB.21 it is the stated closed relation
space. The empty present tuple stays present. External absence
maps to external absence. At any of the already specified
compatible realizations $R\to\mathbb{C}$, these complex-linear
maps are $R$-linear with the given action, and their scalar
extension sends $b\otimes x$ to $b\otimes f(x)$. No new
residue-field realization, elimination of arithmetic nilpotents,
or identification of supported zero with absence is made.
